\chapter{验证定理与最优反馈控制}
\lectureribbon{07}{HJB Equations, DPP and Stochastic Optimal Control II}{Andrzej Święch}
\begin{learninggoals}
\begin{itemize}
  \item 用 Itô 公式证明任何 HJB 光滑解都是成本的下界；
  \item 理解何时 Hamiltonian 的极小点产生真正的最优反馈；
  \item 区分强控制与弱控制表述，以及退化/一致抛物两种 PDE 情形；
  \item 看清验证定理的适用边界。
\end{itemize}
\end{learninggoals}

\section{HJB 候选解怎样“验证”最优性}

令
\[
\mathcal H(t,x,p,M)
=\inf_{a\in A}\left\{
b(t,x,a)\cdot p
+\frac12\Tr[\sigma\sigma^\top(t,x,a)M]
+f(t,x,a)\right\}.
\]
假设 $v\in C^{1,2}$ 满足
\[
\partial_tv+\mathcal H(t,x,Dv,D^2v)=0,
\qquad v(T,x)=g(x).
\]
对任意容许控制 $\alpha$，由于 infimum 不大于把 $a$ 换成 $\alpha_s$ 后的值，
\[
\partial_tv+\mathcal L^{\alpha_s}v+f(s,x,\alpha_s)\ge0.
\]
沿受控状态使用 Itô 公式并取期望，鞅项消失：
\[
\E[g(X_T)]-v(t,x)
=\E\int_t^T(\partial_sv+\mathcal L^{\alpha_s}v)(s,X_s)\dd s.
\]
于是
\[
J(t,x;\alpha)-v(t,x)
=\E\int_t^T
\bigl(\partial_sv+\mathcal L^{\alpha_s}v+f(s,X_s,\alpha_s)\bigr)\dd s
\ge0.
\]

\begin{theorem}[经典验证定理]
若上述光滑性、增长条件与可积性足以合法使用 Itô 公式，则
\[
v(t,x)\le V(t,x).
\]
若存在容许控制 $\alpha^*$ 几乎处处实现 HJB 中的极小值，则等号成立：
\[
v(t,x)=V(t,x)=J(t,x;\alpha^*).
\]
\end{theorem}

\begin{center}
\begin{tikzpicture}[node distance=12mm and 17mm]
  \node[concept] (hjb) {HJB 光滑解\\$v$};
  \node[conceptyellow,right=of hjb] (ineq) {对任意控制\\$v\le J(\alpha)$};
  \node[conceptgreen,right=of ineq] (attain) {若逐点取到极小\\$J(\alpha^*)=v$};
  \draw[flow] (hjb)--node[above,font=\scriptsize\sffamily,text=MutedText]{Itô} (ineq);
  \draw[warmflow] (ineq)--node[above,font=\scriptsize\sffamily,text=MutedText]{等号条件} (attain);
\end{tikzpicture}
\captionof{figure}{验证证明有两步：先得到普遍下界，再找一条使所有不等式变等式的控制。}
\end{center}

\section{从逐点极小化到反馈控制}

若存在可测选择器
\[
\hat a(t,x)\in\argmin_{a\in A}
\left\{\mathcal L^av(t,x)+f(t,x,a)\right\},
\]
自然候选是 Markov 反馈
\[
\alpha_s^*=\hat a(s,X_s^*),
\]
其中闭环状态满足
\[
\dd X_s^*=b(s,X_s^*,\hat a(s,X_s^*))\dd s
+\sigma(s,X_s^*,\hat a(s,X_s^*))\dd W_s.
\]
但“写出 argmin”还不够：必须确保选择器可测、闭环 SDE 有解、控制确实容许且成本可积。

\begin{center}
\begin{tikzpicture}[node distance=13mm and 18mm]
  \node[concept] (x) {观测状态\\$X_s$};
  \node[conceptyellow,right=of x] (policy) {反馈律\\$\hat a(s,X_s)$};
  \node[conceptgreen,right=of policy] (dyn) {闭环动力学\\$\dd X_s$};
  \draw[flow] (x)--(policy);
  \draw[warmflow] (policy)--(dyn);
  \draw[softflow,bend left=32] (dyn) to node[above,font=\scriptsize\sffamily,text=MutedText]{新状态} (x);
\end{tikzpicture}
\captionof{figure}{反馈不是一条预先画好的轨迹，而是状态—动作—新状态的闭环。}
\end{center}

\section{一个随机线性二次例子}

考虑
\[
\dd X_s=(aX_s+b\alpha_s)\dd s+\eta\dd W_s,
\qquad
J=\E\!\left[\int_t^T\frac12(qX_s^2+r\alpha_s^2)\dd s+\frac h2X_T^2\right],
\]
其中 $r>0$。设
\[
V(t,x)=\frac12P(t)x^2+c(t).
\]
HJB 中关于控制的部分是
\[
b\alpha Px+\frac r2\alpha^2,
\]
故
\[
\alpha^*(t,x)=-\frac{bP(t)}r\,x.
\]
比较 $x^2$ 与常数项得到 Riccati 系统
\[
\begin{aligned}
P'(t)+2aP(t)-\frac{b^2}{r}P(t)^2+q&=0,
&&P(T)=h,\\
c'(t)+\frac{\eta^2}{2}P(t)&=0,
&&c(T)=0.
\end{aligned}
\]
噪声不改变线性反馈斜率的形式，却通过 $c(t)$ 提高期望成本。

\section{一致抛物与退化情形}

若存在 $\lambda>0$ 使
\[
\sigma\sigma^\top(t,x,a)\ge\lambda I
\quad\text{对所有 }(t,x,a),
\]
称方程一致抛物。噪声在所有方向平滑，经典正则性更有希望。若 $\sigma\sigma^\top$ 只有低秩，甚至 $\sigma=0$，HJB 是退化抛物或一阶方程，价值函数更容易出现折角。

\begin{tcolorbox}[colback=NightPanelTwo,colframe=GridLine]
\begin{tabularx}{\textwidth}{@{}>{\sffamily\bfseries\color{BlueAccent}}p{0.22\textwidth}YY@{}}
\toprule
 & \textcolor{GreenAccent}{一致抛物} & \textcolor{YellowAccent}{退化/一阶} \\
\midrule
噪声方向 & 覆盖所有状态方向 & 只覆盖部分方向或没有噪声 \\
平滑作用 & 强 & 弱或缺失 \\
经典解 & 条件合适时可期待 & 常常不能期待 \\
稳健语言 & 经典解与黏性解 & 黏性解通常不可缺少 \\
\bottomrule
\end{tabularx}
\end{tcolorbox}

\section{强表述与弱表述}

\begin{definition}[强控制问题]
概率空间、Brown 运动和滤过预先固定；控制在这套随机基底上选择，状态是给定 Brown 运动的强解。
\end{definition}

\begin{definition}[弱控制问题]
优化时允许概率空间、Brown 运动乃至状态联合分布一并变化，只要求它们满足相应的鞅问题或 SDE 分布约束。
\end{definition}

强表述便于构造路径上的控制；弱表述在扩散系数受控、强解不唯一或需要紧性方法时更灵活。两者是否有相同值，需要额外条件，不能想当然。

\begin{pitfall}[验证定理不负责制造光滑性]
验证定理的逻辑是“若有足够光滑的候选解，则可证明它就是值函数”。它不证明 HJB 必有经典解。现实中值函数常只连续，此时要用黏性解与比较原理替代逐点二阶导数。
\end{pitfall}

\begin{takeaway}
\begin{enumerate}
  \item HJB 光滑解经 Itô 公式给出所有控制成本的下界；
  \item 逐点实现 Hamiltonian 极小值且闭环可解时，反馈控制最优；
  \item 可测选择、闭环适定性和可积性是 argmin 背后的三道门；
  \item 退化扩散会削弱平滑性，推动理论转向黏性解。
\end{enumerate}
\end{takeaway}

\begin{exercisebox}
\begin{enumerate}
  \item 在验证推导中逐行指出哪一步使用了 HJB，哪一步使用了终端条件。
  \item 完成标量 Riccati 例子的系数比较，并说明 $r$ 变大时反馈强度如何变化。
  \item 举出一个逐点 argmin 存在但闭环 SDE 可能不适定的原因。
\end{enumerate}
\end{exercisebox}
